% BHU PYQ -- Home Science: Textile Surface Design and Basic Stitching (2024-25 NEP)
\documentclass[11pt,a4paper]{article}
\usepackage[a4paper,top=2.5cm,bottom=2.5cm,left=2.8cm,right=2.8cm,headheight=14pt]{geometry}
\usepackage{amsmath,amssymb}
\usepackage[T1]{fontenc}
\usepackage[utf8]{inputenc}
\usepackage{lmodern,microtype,fancyhdr,enumitem,hyperref,lastpage,parskip}

\hypersetup{colorlinks=true,urlcolor=black,linkcolor=black,pdftitle={HSCSE21: Textile Surface Design and Basic Stitching -- BHU NEP 2024-25}}

\pagestyle{fancy}\fancyhf{}
\renewcommand{\headrulewidth}{0.4pt}\renewcommand{\footrulewidth}{0.4pt}
\fancyhead[L]{\small \textit{Banaras Hindu University}}
\fancyhead[R]{\small \textit{HSCSE21: Textile Surface Design (2024-25)}}
\fancyfoot[C]{\small Page \thepage\ of \pageref{LastPage}}

\newlist{parts}{enumerate}{2}
\setlist[parts,1]{label=(\alph*),leftmargin=*,itemsep=4pt,topsep=3pt}
\setlist[parts,2]{label=(\roman*),leftmargin=*,itemsep=2pt,topsep=2pt}

\newcommand{\pts}[1]{\hfill \textit{[#1]}}

\begin{document}

\begin{center}
  {\large\bfseries BANARAS HINDU UNIVERSITY}\\[2pt]
  {\normalsize B. A. / B. Sc. (Hons.) Semester II Examination, 2024-25 (NEP)}\\[6pt]
  \rule{0.8\linewidth}{0.5pt}\\[6pt]
  {\large\bfseries HOME SCIENCE}\\[4pt]
  {\large\bfseries Paper Code: HSCSE-21 : Textile Surface Design and Basic Stitching}\\[4pt]
  {\normalsize Time: 2:30 Hours \hfill Full Marks: 70}\\[2pt]
  {\small REF NO. CE/Conf./D-II/108}
\end{center}
\rule{\linewidth}{0.4pt}
\smallskip

\noindent \textit{Instructions: This question paper comprises 3 Sections A, B and C. See the instructions of each section and attempt necessary questions. Write your Roll No.\ at the top immediately on receipt of this question paper.}

\bigskip

\section*{SECTION -- A}
\noindent \textit{Note: Attempt any \textbf{three} of the following (each within 250 words).}\pts{3 $\times$ 10 = 30}

\begin{enumerate}[label=\textbf{\arabic*.},leftmargin=*,itemsep=10pt]

  \item What are the ``Elements of Design''? Explain with illustrations.\pts{10}

  \item Explain any two methods of surface designing through dyes and prints.\pts{10}

  \item Discuss any five types of Embroidery stitches.\pts{10}

  \item What are the different parts of a sewing machine? Explain them.\pts{10}

\end{enumerate}

\bigskip

\section*{SECTION -- B}
\noindent \textit{Note: Attempt any \textbf{four} of the following (each within 100 words).}\pts{4 $\times$ 5 = 20}

\begin{enumerate}[label=\textbf{\arabic*.},leftmargin=*,start=5,itemsep=8pt]

  \item Discuss the concept and importance of textile surface designing.\pts{5}

  \item Explain the uses of Buttonhole stitch through a diagram.\pts{5}

  \item Discuss about any two methods of Tie \& Dye technique.\pts{5}

  \item What do you understand by Applique work? Explain.\pts{5}

  \item Differentiate between stitch and seam. Also write about the French seam.\pts{5}

\end{enumerate}

\bigskip

\section*{SECTION -- C}
\noindent \textit{Note: Write about the following within 50 words each:}\pts{10 $\times$ 2 = 20}

\begin{enumerate}[label=\textbf{10.}(\roman*),leftmargin=*,itemsep=6pt]

  \item Rhythm in design\pts{2}
  \item Lahariya\pts{2}
  \item Slot Seam\pts{2}
  \item Uses of pressure foot\pts{2}
  \item Fasteners\pts{2}
  \item Batik\pts{2}
  \item Box Pleats\pts{2}
  \item Pin Tucks\pts{2}
  \item Yokes\pts{2}
  \item Symmetrical Balance\pts{2}

\end{enumerate}

\end{document}
